% Compile with pdfLaTeX (two passes), or upload directly to Overleaf.
\documentclass[11pt,a4paper]{article}
\usepackage[T1]{fontenc}
\usepackage[utf8]{inputenc}
\usepackage[margin=2.5cm]{geometry}
\usepackage{lmodern}
\usepackage{amsmath}
\usepackage{booktabs,longtable,tabularx,array}
\usepackage{xurl}
\usepackage[hidelinks]{hyperref}
\usepackage{listings}
\usepackage{xcolor}
\usepackage{caption}
\renewcommand{\contentsname}{Daftar Isi}
\renewcommand{\abstractname}{Ringkasan}
\renewcommand{\tablename}{Tabel}
\renewcommand{\refname}{Daftar Pustaka}
\newcommand{\kode}[1]{\texttt{\detokenize{#1}}}
\newcolumntype{L}[1]{>{\raggedright\arraybackslash}p{#1}}
\lstset{basicstyle=\ttfamily\small,breaklines=true,
 frame=single,columns=fullflexible,keepspaces=true,
 showstringspaces=false,backgroundcolor=\color{gray!5}}
\setlength{\parskip}{0.45em}
\setlength{\parindent}{0pt}
\title{Laporan Kemajuan Mingguan\\
 Pemetaan Fungsi Program dan Tuning Attention Patch Size\\
 pada Point Transformer V3}
\author{Nama: Nicholas Shane Pangihutan Siahaan\\
 NIM : 23/520590/TK/57399\\
 Program studi: Teknologi Informasi}
\date{16 September 2026}

\begin{document}
\maketitle
\begin{center}
\textbf{Proyek:} Transformer-based Models for Geometric Deep Learning
on Point-cloud Data\\
\textbf{Periode eksperimen:} 09--15 September 2026
\end{center}

% \begin{abstract}
% Laporan ini menyajikan pemetaan fungsi utama pipeline klasifikasi
% Point Transformer V3 (PTv3) pada framework Pointcept serta eksperimen
% tuning ukuran patch attention. Eksperimen menggunakan subset stratified
% ModelNet40 yang terdiri atas 798 objek training dan 202 objek validation.
% Seluruh konfigurasi menggunakan 1.024 titik per objek, seed 42, dan
% 100 epoch. Ukuran patch encoder 128, 256, dan 512 dibandingkan dengan
% baseline 1024. Baseline mencapai validation overall accuracy tertinggi
% sebesar 75,25\%, sedangkan patch 256 mencapai 74,75\% dengan mean class
% accuracy sedikit lebih tinggi. Waktu training seluruh konfigurasi berada
% di sekitar 47--48 menit. Hasil satu seed ini belum menunjukkan peningkatan
% overall accuracy atau percepatan yang berarti akibat pengecilan patch.
% \end{abstract}
% \tableofcontents
% \newpage

% \section{Latar Belakang dan Tujuan}
% Point cloud memiliki susunan titik yang tidak mengikuti grid reguler.
% Model perlu mengolah fitur dan hubungan geometris tanpa bergantung pada
% urutan input yang arbitrer. PTv3 menggunakan serialized neighborhood
% mapping untuk membentuk kelompok attention secara efisien, sebagai
% alternatif pencarian neighborhood presisi berbasis KNN pada pipeline
% sebelumnya~\cite{ptv3}.
\section{Tujuan}
Tujuan kegiatan minggu ini adalah:
\begin{enumerate}
\item Memetakan fungsi program, input, output, dan hubungan pemanggilan
      dalam pipeline PTv3 untuk klasifikasi ModelNet40.
\item Memahami serialization, proyeksi Query--Key--Value (QKV), patch
      attention, hierarchical pooling, dan update parameter.
\item Menguji pengaruh \kode{enc_patch_size} dengan kondisi eksperimen
      lainnya tetap.
\item Mendokumentasikan hasil validation dan durasi training sebagai
      dasar eksperimen lanjutan.
\end{enumerate}
Kegiatan ini merupakan \textit{hyperparameter tuning}; model dilatih
dari awal. Belum dilakukan fine-tuning dari pretrained checkpoint maupun
modifikasi rumus QKV.

\section{Pemetaan Struktur dan Fungsi Program}
\subsection{Struktur repository}
% \begin{tabularx}{\linewidth}{lX}
% \toprule
% Bagian & Fungsi utama \\
% \midrule
\kode{configs/} : Parameter model, dataset, optimizer, scheduler, dan hooks. \\
\kode{pointcept/datasets/} : Membaca data, sampling, dan transformasi. \\
\kode{pointcept/models/} : Backbone PTv3, classification head, dan loss. \\
\kode{pointcept/engines/} : Training, validation, testing, dan hooks. \\
\kode{pointcept/utils/} : Config parser, logging, registry, dan optimizer. \\
\kode{libs/} : Operasi tambahan point cloud berbasis C++/CUDA. \\
\kode{tools/} : Entry point Python dan program pembentukan split dataset. \\
\kode{scripts/} : Peluncur shell untuk training dan testing. \\
\kode{data/} : Dataset dan daftar split; pada workspace ini berupa symlink. \\
\kode{exp/} : Log, checkpoint, TensorBoard, config, dan snapshot kode. \\
% \bottomrule
% \end{tabularx}

\subsection{Hubungan pemanggilan}
\begin{lstlisting}
scripts/train.sh
  -> tools/train.py: main()
  -> default_config_parser() -> launch() -> main_worker()
  -> Trainer.__init__()
     -> model, dataset, DataLoader, optimizer, scheduler
  -> Trainer.train()
     -> dataset.__getitem__() -> transforms -> batch
     -> Trainer.run_step()
        -> DefaultClassifier.forward()
           -> PointTransformerV3.forward()
           -> mean aggregation -> classification head -> loss
        -> backward -> optimizer.step() -> scheduler.step()
     -> validation -> save checkpoint
  -> exp/modelnet40/<nama-eksperimen>/
\end{lstlisting}
\kode{scripts/train.sh} menyalin source code ke \kode{exp/.../code/}
dan menjalankan snapshot tersebut. Perubahan source setelah training
dimulai tidak otomatis mengubah proses yang sedang aktif.

\subsection{Peta fungsi utama}
Path yang diawali \kode{datasets/}, \kode{models/}, dan \kode{engines/}
pada tabel berikut berada di dalam folder \kode{pointcept/}.
\small
\renewcommand{\arraystretch}{1.15}
\renewcommand{\kode}[1]{\nolinkurl{#1}}
\setlength{\tabcolsep}{4pt}
\begin{longtable}{@{} L{0.34\linewidth}L{0.27\linewidth}L{0.29\linewidth}@{}}
\caption{Peta fungsi utama pipeline klasifikasi PTv3.}\label{tab:functions}\\
\toprule
File / fungsi & Input dan proses & Output / efek \\
\midrule
\endfirsthead
\toprule
File / fungsi & Input dan proses & Output / efek \\
\midrule
\endhead
\bottomrule
\endfoot
\kode{tools/train.py: main()} & Membaca argumen CLI dan config; meluncurkan worker. & Proses training GPU. \\
\kode{main_worker(cfg)} & Setup run; membangun trainer. & Memanggil \kode{trainer.train()}. \\
\kode{engines/defaults.py: default_config_parser()} & Menggabungkan config dan override; menetapkan dataset loop. & Config run dan salinannya. \\
\kode{default_setup()} & Menetapkan seed serta batch/worker per GPU. & Environment eksperimen. \\
\kode{datasets/modelnet.py: get_data_list()} & Membaca \kode{modelnet40_<split>.txt}. & Daftar nama objek. \\
\kode{ModelNetDataset.__init__()} & Root, split, jumlah titik, dan transform; membaca/membuat cache. & Dataset siap digunakan. \\
\kode{get_data(idx)} & Cache atau file XYZ+normal; sampling titik; pembentukan label. & Dictionary koordinat, normal, kategori. \\
\kode{__getitem__(idx)} & Memilih jalur persiapan train atau test. & Satu sampel. \\
\kode{__len__()} & Jumlah objek dikalikan dataset loop. & Panjang dataset. \\
\kode{prepare_train_data(idx)} & Mengambil data dan menerapkan transform. & Input train/validation. \\
\kode{prepare_test_data(idx)} & Mengambil data dan membuat augmentasi test. & Voting list dan label. \\
\kode{datasets/transform.py} & Normalisasi, scale, shift, grid sampling, konversi tensor. & Koordinat, grid coordinate, fitur dan label terpilih. \\
\kode{models/default.py: DefaultClassifier.forward()} & Menjalankan backbone, mean aggregation per objek, head dan loss. & Loss saat train; logits dan loss sesuai mode evaluasi. \\
\kode{models/utils/structure.py: Point.serialization()} & Grid coordinates, batch, dan metode ordering. & Serialized code, order, dan inverse. \\
\kode{models/utils/serialization/default.py: encode()} & Memilih Z-order atau Hilbert; varian trans menukar x dan y. & Kode spasial integer. \\
\kode{PointTransformerV3.forward()} & Serialization, sparsify, embedding dan encoder. & Fitur backbone dalam struktur Point. \\
\kode{Embedding.forward()} & Fitur awal dan struktur sparse point cloud. & Fitur embedding awal. \\
\kode{Block.forward()} & CPE, normalisasi, attention, MLP dan residual. & Fitur hasil transformer block. \\
\kode{SerializedAttention.get_padding_and_inverse()} & Offset tiap objek dan patch size; membentuk padding. & pad, unpad dan cu\_seqlens. \\
\kode{SerializedAttention.forward()} & Fitur terurut; proyeksi QKV, attention patch dan proyeksi output. & Fitur hasil attention. \\
\kode{MLP.forward()} & Linear, GELU, linear dan dropout. & Fitur hasil feed-forward network. \\
\kode{SerializedPooling.forward()} & Pengelompokan hierarkis dan agregasi fitur. & Lebih sedikit titik dengan fitur stage berikutnya. \\
\kode{SerializedUnpooling.forward()} & Upsampling dan skip features pada decoder. & Resolusi lebih tinggi; tidak digunakan saat \kode{enc_mode=True}. \\
\kode{engines/train.py: build_model()} & Membentuk model dari config dan menempatkannya di GPU. & Model trainable. \\
\kode{build_train_loader()}, \kode{build_val_loader()} & Membentuk dataset, sampler, dan DataLoader. & Iterator batch. \\
\kode{build_optimizer()}, \kode{build_scheduler()} & Parameter model dan config update bobot. & Optimizer dan scheduler. \\
\kode{Trainer.train()} & Iterasi epoch/batch dan pemanggilan hooks. & Proses training lengkap. \\
\kode{Trainer.run_step()} & Pemindahan batch ke CUDA, forward, backward dan update. & Bobot baru dan loss batch. \\
\kode{hooks/evaluator.py: ClsEvaluator} & Prediksi pada validation loader. & allAcc, mAcc, F1 dan statistik kelas. \\
\kode{hooks/misc.py: CheckpointSaver} & Epoch dan skor validation. & \kode{model_best.pth}, \kode{model_last.pth}. \\
\kode{PreciseEvaluator.after_train()} & Checkpoint dan test config. & Testing otomatis setelah training. \\
\kode{engines/test.py: ClsVotingTester.test()} & Model dan augmentasi objek test. & Hasil voting evaluation. \\
\end{longtable}
\normalsize
Komponen backbone pada tabel terdapat di
\path{pointcept/models/point_transformer_v3/point_transformer_v3m1_base.py}.

\subsection{QKV dan serialized attention}
Untuk fitur $X$, proyeksi attention secara konseptual ditulis sebagai
\begin{equation}
Q=XW_Q+b_Q,\quad K=XW_K+b_K,\quad V=XW_V+b_V.
\end{equation}
Pada implementasi ini, ketiga proyeksi dikemas dalam satu layer linear
\kode{Linear(channels, 3 * channels)}. Bobot proyeksi dipelajari selama
training. Attention per head dalam satu patch secara konseptual adalah
\begin{equation}
\operatorname{Attention}(Q,K,V)=
\operatorname{softmax}\left(\frac{QK^{\mathsf T}}{\sqrt{d_k}}\right)V.
\end{equation}
Fitur titik diurutkan menggunakan kode spasial sebelum dibagi menjadi
patch. Flash Attention digunakan untuk menghitung attention pada pipeline
aktif. Koordinat juga memengaruhi grid, ordering dan conditional positional
encoding (CPE). Pada konfigurasi ini relative positional encoding eksplisit
dinonaktifkan, yaitu \kode{enable_rpe=False}.

\section{Eksperimen}
\subsection{Dataset}
Training set ModelNet40 terdiri atas 9.843 objek. Program
\path{tools/create_modelnet_subset.py} memilih sekitar 10\% objek dari
setiap kelas secara acak dengan seed 42 dan pembulatan ke atas,
menghasilkan 1.000 objek (sekitar 10,16\% dari training set asli).
Program \path{tools/split_modelnet_train_val.py} membagi subset tersebut
secara stratified sekitar 80:20 dengan pembulatan per kelas.
\begin{table}[htbp]
\centering
\caption{Pembagian dataset yang tetap untuk seluruh eksperimen.}
\begin{tabular}{lrl}
\toprule
Bagian & Objek & Nama split \\
\midrule
Training & 798 & \kode{train10_train} \\
Validation & 202 & \kode{train10_val} \\
Official test & 2.468 & \kode{test} \\
\bottomrule
\end{tabular}
\end{table}
Seluruh 40 kelas terwakili pada training dan validation. Kedua split tidak
beririsan dan gabungannya sama dengan subset awal. Official test juga
terpisah dari keduanya. Stratified random sampling mempertahankan
distribusi kelas secara mendekati, tetapi tidak menjamin seluruh variasi
geometri intra-kelas terwakili.

Pada level titik, farthest point sampling (FPS) mereduksi setiap objek
menjadi 1.024 titik sebelum transformasi grid. Ini berbeda dari pemilihan
subset antarobjek. Grid sampling selanjutnya dapat mengurangi jumlah
titik aktual yang masuk ke model.

\subsection{Environment dan konfigurasi tetap}
Eksperimen dijalankan pada laptop dengan GPU NVIDIA GeForce RTX 3060
Laptop, Sistem Operasi Linux, Distro Ubuntu 22.04 dan VRAM 6 GB. Environment yang digunakan adalah
\kode{pointcept-torch2.5.0-cu12.4}: Python 3.10, PyTorch 2.5.0 dengan
CUDA 12.4, spconv 2.3.8, Flash Attention 2.7.4.post1, dan torch-scatter
2.1.2. Versi ini berasal dari pemeriksaan environment proyek.

\begin{tabularx}{\linewidth}{lX}
\toprule
Parameter & Nilai \\
\midrule
Backbone / task & PT-v3m1 / DefaultClassifier, 40 kelas \\
Mode & Encoder-only, \kode{enc_mode=True} \\
Jumlah parameter & 38.842.792 berdasarkan log \\
Input features & Koordinat XYZ dan normal XYZ, total 6 channel \\
Seed model / split & 42 / 42 \\
Epoch / eval\_epoch & 100 / 100; train.loop = 1 \\
Batch train / validation & 4 / 8 \\
Worker & 2 \\
Optimizer & AdamW; LR dasar 0,001; weight decay 0,01 \\
LR kelompok block & 0,0001; sesuai \kode{param_dicts} \\
Scheduler & OneCycleLR; max LR kelompok 0,001 dan 0,0001 \\
Loss & Cross Entropy + Lovasz multiclass, masing-masing bobot 1 \\
Ordering & z, z-trans, hilbert, hilbert-trans; shuffle aktif \\
Regularisasi & drop\_path = 0,3; attention dropout = 0 \\
Flash / AMP & Flash Attention aktif; AMP global nonaktif \\
\bottomrule
\end{tabularx}
Meskipun AMP global nonaktif, implementasi Flash Attention mengonversi
QKV ke dtype yang didukung kernel di dalam modul attention.

\subsection{Variabel tuning dan kontrol}
Variabel yang diubah adalah \kode{enc_patch_size}, diterapkan sama pada
kelima stage encoder. Nilai kandidat adalah 128, 256, dan 512;
baseline menggunakan 1024. Semua model dimulai dari random initialization
tanpa memuat checkpoint baseline. Dataset, optimizer, scheduler, seed,
jumlah epoch, ordering, dan parameter model lainnya tetap.

Untuk mencegah testing otomatis selama tuning, config induk
\kode{cls-ptv3-v1m1-0-laptop-tuning.py} menghilangkan hook
\kode{PreciseEvaluator}. Validation dan checkpoint tetap berjalan.
Baseline sebelumnya memang menjalankan official test setelah training,
tetapi angka official test tidak digunakan untuk memilih patch size.

\subsection{Metrik dan durasi}
Metrik pemilihan model adalah overall accuracy pada validation:
\begin{equation}
\mathrm{allAcc}=\frac{\sum_{c=1}^{40}n_{c,\mathrm{correct}}}
{\sum_{c=1}^{40}n_c},\qquad
\mathrm{mAcc}=\frac{1}{40}\sum_{c=1}^{40}
\frac{n_{c,\mathrm{correct}}}{n_c}.
\end{equation}
Macro F1 merupakan rata-rata F1 per kelas. Setiap model dipilih pada epoch
dengan validation allAcc tertinggi; jika skor sama, tabel menggunakan
kemunculan pertama. mAcc dan macro F1 dilaporkan pada epoch terpilih,
bukan maksimum terpisah sepanjang training.

Durasi dihitung dari timestamp \kode{Start Training} sampai hasil
validation epoch ke-100. Durasi mencakup training dan validation berkala,
tetapi tidak mencakup persiapan dataset, voting test, atau penulisan
checkpoint terakhir setelah validation. Durasi bukan benchmark inference.

\section{Hasil Eksperimen}
Semua run selesai 100 epoch tanpa error yang tercatat di log. Setiap run
menghasilkan \kode{model_best.pth} dan \kode{model_last.pth}.
\begin{table}[htbp]
\centering
\caption{Hasil pada validation set. Seluruh skor dalam persen.}
\label{tab:results}
\begin{tabular}{rrrrrl}
\toprule
Patch & allAcc & mAcc & Macro F1 & Epoch & Durasi \\
\midrule
1024 & \textbf{75,25} & 60,81 & 57,50 & 73 & 47 m 36 d \\
128  & 73,76 & 57,39 & 53,08 & 83 & 47 m 15 d \\
256  & 74,75 & \textbf{61,14} & 57,45 & 80 & 47 m 55 d \\
512  & 74,26 & 60,75 & \textbf{58,90} & 95 & 48 m 25 d \\
\bottomrule
\end{tabular}
\end{table}
Pada tabel, m dan d berarti menit dan detik; waktu dibulatkan ke detik
terdekat. Skor dibulatkan dari log. Dengan 202 objek validation, jumlah
prediksi benar pada epoch terbaik adalah 152, 149, 151, dan 150 untuk
patch 1024, 128, 256, dan 512 secara berurutan.

\section{Analisis dan Keterbatasan}
\subsection{Akurasi}
Baseline patch 1024 memberikan allAcc tertinggi. Patch 128, 256, dan
512 lebih rendah sekitar 1,49; 0,50; dan 0,99 poin persentase.
Patch 256 menjadi kandidat tuning paling dekat dengan baseline, dan
mAcc-nya sedikit lebih tinggi. Patch 512 memiliki macro F1 tertinggi,
menunjukkan bahwa peringkat konfigurasi dapat berbeda menurut metrik.
Pemilihan utama tetap mengikuti allAcc yang telah ditetapkan.

Satu prediksi benar tambahan pada validation set mengubah allAcc sebesar
$100/202\approx0,495$ poin persentase. Selisih baseline dan patch 256
hanya satu objek. Oleh karena itu, hasil ini tidak cukup untuk membuktikan
superioritas statistik konfigurasi tertentu.

\subsection{Efisiensi}
Durasi seluruh run berada pada rentang 47 menit 15 detik sampai
48 menit 25 detik. Patch 128 sekitar 21 detik lebih cepat dari baseline,
sedangkan patch 256 dan 512 sedikit lebih lambat. Pengecilan patch belum
memberikan penghematan waktu training yang berarti pada pengukuran ini.

Hipotesis penjelasannya adalah bahwa input hanya 1.024 titik dan pooling
mengurangi jumlah titik pada stage berikutnya; sebagian stage mungkin
memiliki jumlah titik di bawah batas patch. Selain itu, durasi total
dipengaruhi sparse convolution, DataLoader, validation, dan penyimpanan
checkpoint, bukan attention saja. Hipotesis ini memerlukan profiling
untuk dibuktikan. Kompleksitas attention konvensional per patch
$O(P^2)$ tidak langsung menentukan keseluruhan waktu pipeline Flash
Attention pada eksperimen ini.

\subsection{Batas validitas}
\begin{enumerate}
\item Seluruh konfigurasi baru diuji dengan satu seed training.
\item Validation set kecil, khususnya kelas minoritas yang hanya memiliki
      satu sampai beberapa objek; mAcc dapat berfluktuasi.
\item Subset mempertahankan distribusi kelas, bukan jaminan keragaman
      geometri lengkap dari dataset asli.
\item Waktu diukur dari log pada sesi berbeda, tanpa kontrol ketat suhu,
      daya GPU, dan proses latar belakang.
\item Peak VRAM, latency inference, dan profiler attention belum diukur.
\item Kurva learning dan confusion matrix belum dianalisis dalam laporan
      ini. Tidak disajikan grafik sintetis atau data yang belum diekstrak.
\end{enumerate}
Angka hasil tidak dapat langsung dibandingkan dengan benchmark paper PTv3
karena dataset, subset, tugas, dan protokol eksperimennya berbeda.

\section{Kesimpulan dan Rencana Selanjutnya}
Pipeline program telah dipetakan dari peluncur training sampai dataset,
backbone, classification head, loss, update bobot, validation, dan
checkpoint. Eksperimen ini berhasil menguji satu hyperparameter dengan
kondisi lain tetap dan tanpa official test otomatis pada run tuning.

Pengecilan patch belum meningkatkan validation allAcc maupun memberi
percepatan total yang berarti. Baseline 1024 tetap menjadi kontrol utama.
Patch 256 disimpan sebagai kandidat konfirmasi karena hasilnya dekat
dengan baseline dan mAcc sedikit lebih tinggi. Temuan ini merupakan
hasil sementara, bukan klaim bahwa patch 1024 selalu optimal.

Rencana selanjutnya:
\begin{enumerate}
\item Menguji ordering Z, Hilbert, dan kombinasi Z+Hilbert dengan patch
      1024 serta parameter lain tetap.
\item Mengulangi baseline dan konfigurasi menjanjikan dengan seed 123
      dan 501; melaporkan mean dan standard deviation.
\item Mengekstrak kurva loss/validation dan memprofilkan komponen waktu
      serta penggunaan VRAM.
\item Melakukan tuning learning rate/regularisasi setelah ablation utama.
\item Melakukan fine-tuning hanya setelah sumber pretrained checkpoint
      dan protokol pembanding ditetapkan.
\item Menjalankan official test pada model terpilih dengan aturan voting
      yang sama dan telah ditetapkan; tidak memilih hasil voting berdasarkan
      skor test tertinggi saja.
\end{enumerate}

% \appendix
% \section{Repositori dan Log Eksperimen}
% Root repository: \path{~/Pointcept}. Commit Git sebagai referensi
% source upstream:
% \path{9f37497e4f3005c90bbbe7221b86439c29d60611}.
% Config dan program split tambahan merupakan file lokal; commit upstream
% saja tidak cukup untuk mereproduksi seluruh eksperimen.
% 
% \begin{longtable}{L{0.19\linewidth}L{0.73\linewidth}}
% \toprule
% Run & Log hasil (relatif terhadap root) \\
% \midrule
% Baseline & \path{exp/modelnet40/baseline-10pct-split/train.log} \\
% Patch 128 & \path{exp/modelnet40/tuning-patch128-s42/train.log} \\
% Patch 256 & \path{exp/modelnet40/tuning-patch256-s42/train.log} \\
% Patch 512 & \path{exp/modelnet40/tuning-patch512-s42/train.log} \\
% \bottomrule
% \end{longtable}
% Setiap folder run menyimpan \kode{config.py}, \kode{code/}, event
% TensorBoard, dan folder \kode{model/}. Salinan config pada folder run
% merupakan sumber pengaturan aktual; snapshot kode merupakan sumber
% implementasi aktual saat eksperimen dijalankan.
% 
% \begin{lstlisting}[language=bash]
% conda activate pointcept-torch2.5.0-cu12.4
% sh scripts/train.sh -g 1 -d modelnet40 \
%   -c cls-ptv3-v1m1-0-laptop-patch128 \
%   -n tuning-patch128-s42
% \end{lstlisting}
% Untuk patch 256 dan 512, ganti angka 128 pada config dan nama run.
% Gunakan nama run baru jika mengulang eksperimen agar hasil lama tidak
% tertimpa. Jangan menjalankan beberapa training bersamaan pada GPU laptop.
% 
% \section{Informasi Hasil Official Test Baseline}
% Log baseline mencatat evaluasi official test sesudah training dengan
% 10 voting evaluation. Record terbaik yang dilaporkan program adalah
% allAcc 74,96\% dan mAcc 61,20\%; evaluasi terakhir menghasilkan
% allAcc 74,55\% dan mAcc 60,61\%. Angka record terbaik ini adalah statistik
% yang dicatat program, bukan hasil protokol final yang memilih voting
% tanpa melihat label test. Angka tersebut tidak digunakan dalam tabel
% ablation patch maupun untuk memilih konfigurasi tuning.
% 
\begin{thebibliography}{9}
\bibitem{ptv3}
X. Wu, L. Jiang, P.-S. Wang, Z. Liu, X. Liu, Y. Qiao, W. Ouyang,
T. He, dan H. Zhao.
\textit{Point Transformer V3: Simpler, Faster, Stronger}.
CVPR, 2024. \url{https://arxiv.org/abs/2312.10035}.
\bibitem{pointcept}
Pointcept contributors. \textit{Pointcept: A Codebase for Point Cloud
Perception}. Repository source code.
\url{https://github.com/Pointcept/Pointcept}.
\bibitem{local}
Artefak eksperimen lokal. Config, snapshot source, checkpoint, dan log
baseline serta tuning patch size, 13--15 September 2026.
Path artefak dicantumkan pada Lampiran A.
\end{thebibliography}
\end{document}
